\documentclass[conference]{IEEEtran}
\IEEEoverridecommandlockouts
% The preceding line is only needed to identify funding in the first footnote. If that is unneeded, please comment it out.
\usepackage{cite}
\usepackage{amsmath,amssymb,amsfonts}
\usepackage{algorithmic}
\usepackage{graphicx}
\usepackage{textcomp}
\usepackage{xcolor}
\def\BibTeX{{\rm B\kern-.05em{\sc i\kern-.025em b}\kern-.08em
    T\kern-.1667em\lower.7ex\hbox{E}\kern-.125emX}}
\begin{document}

\title{Kinematic Decoupling of Fission Recycling in Kilonovae: A Delayed Neutron Release Mechanism for Third r-Process Peak Shifts and Anisotropic Curtaining}\author{Anonymous Research Plan Author}\maketitle

\begin{abstract}
Multi-messenger observations and nucleosynthetic networks confirm that binary neutron star mergers produce lanthanide-rich ejecta and synthesize heavy r-process elements through fission recycling. This proposal advances a mathematical theory formalizing the kinematic decoupling of fission recycling, positing that delayed neutron release from beta-decaying and isomer-trapped superheavy nuclei modifies the third r-process peak abundance. The central mechanism asserts that when the effective fission timescale exceeds the ejecta expansion timescale, the violation of instantaneous equilibrium induces a distinct temporal evolution of actinide spectral lines. This temporal signature is proposed to be uniquely sensitive to fission kinetics, establishing a formal boundary against alternative explanations driven purely by geometric or magnetic ejecta asymmetries.

The theoretical framework requires a rigorous discrimination proof to demonstrate the non-isomorphism between delayed fission spectral evolution and static geometric or magnetic asymmetries. The planned contribution relies on deriving an analytical relationship between the effective fission timescale, the expansion timescale, and abundance yields. Because the mathematical definitions for these timescales and the operational metrics for spectral evolution remain unresolved, the design proceeds as an unverified proposal. The expected outcome is a formal proof obligation that either supports the proposed mechanism within declared boundaries or invalidates it if the fission timescale is strictly shorter than the expansion timescale.
\end{abstract}
\begin{IEEEkeywords}
kilonovae, r-process, fission recycling, delayed neutron release, third r-process peak, lanthanide curtaining, actinide spectral lines, kinematic delay
\end{IEEEkeywords}\section{Introduction}
\subsection{Background}
Binary neutron star mergers are established as the dominant astrophysical site for rapid neutron-capture (r-process) nucleosynthesis, synthesizing elements heavier than iron under extreme neutron-rich conditions. Multi-messenger observations of GW170817 and its kilonova counterpart AT2017gfo confirmed the presence of lanthanide-rich ejecta, characterized by high opacity and redshifted near-infrared emission. The observed light curves and spectral features align with models of radioactive decay powering the transient, where the electron fraction (\$Y\_e\$) dictates the nucleosynthetic pathway. Low \$Y\_e\$ trajectories (\$Y\_e \textbackslash{}lesssim 0.2\$) drive the flow toward the third r-process peak and actinides, while higher \$Y\_e\$ components produce lighter r-process elements. This compositional stratification is critical for interpreting the electromagnetic signatures of merger ejecta.~\cite{cite_07d6df5d6978a1fc,cite_30c88fe46f23448d,cite_612290131318c0d4}

\subsection{Research Gap}
Despite this consensus, a significant theoretical gap remains in linking the timing of fission recycling to observable spectral evolution. Current models often assume instantaneous fission equilibrium, treating fission neutron release as prompt relative to ejecta expansion. However, beta-delayed and isomer-delayed fission introduce a kinematic delay, potentially decoupling the late-stage r-process from the prompt outflow. This delay modifies the third r-process peak freeze-out abundance and induces viewing-angle-dependent lanthanide curtaining. Crucially, the resulting temporal evolution of actinide spectral lines is hypothesized to be distinct from patterns generated by static geometric or magnetic asymmetries. Existing literature lacks a formal framework to discriminate between kinematic fission delays and purely geometric confounders in kilonova spectra, leaving the specific spectral signatures of delayed fission unverified.~\cite{cite_43194f3b6113c2ab,cite_b668987a568d271f}

\subsection{Proposed Contribution}
This research plan proposes a mathematical theory of kinematic decoupling in fission recycling. The central contribution is a formal unification of third r-process peak shifts and anisotropic lanthanide curtaining through a delayed neutron release mechanism. We introduce a static axiomatic framework where the effective fission timescale (\$t\_\{fission\}\$) competes with the ejecta expansion timescale (\$t\_\{expansion\}\$). The proposal asserts that when \$t\_\{fission\} \textbackslash{}ge t\_\{expansion\}\$, the violation of instantaneous equilibrium produces a unique temporal evolution of actinide spectral lines. This mechanism is posited to be non-isomorphic to spectral signatures arising from geometric asymmetries or magnetic field configurations alone. The work establishes the theoretical basis for distinguishing kinematic fission delays from static ejecta morphologies.

\subsection{Design Assumptions}
The theory rests on two core design assumptions. First, the effective fission timescale, governed by beta-decay half-lives and nuclear isomeric states, is comparable to or longer than the ejecta expansion timescale in high-entropy, low-density trajectories. Second, actinide spectral line evolution is uniquely sensitive to this fission timescale and cannot be mimicked by variations in ejecta geometry or magnetic field topology. These assumptions define the boundary conditions for the proposed mechanism. The failure of the mechanism is explicitly scoped to regimes where \$t\_\{fission\} < t\_\{expansion\}\$ or where geometric/magnetic asymmetries perfectly replicate the predicted spectral shifts without invoking delayed fission. Formal definitions for the timescales and spectral metrics are required to advance the proof obligations.

\section{Background, Survey, and Research Gap}
\subsection{Established Site and Nucleosynthetic Benchmark}
The rapid neutron-capture process (r-process) synthesizes elements heavier than iron in transient, high-energy environments where neutron capture outpaces beta decay. Binary neutron star mergers are the dominant astrophysical site for the main r-process, confirmed by the multi-messenger detection of GW170817 and its kilonova counterpart AT2017gfo. The electromagnetic emission of these transients is powered by the radioactive decay of neutron-rich ejecta, producing a characteristic two-component structure: a fast, lanthanide-poor blue component and a slower, lanthanide-rich red component. This observational anchor establishes the physical regime---extreme neutron densities and rapid expansion---within which fission recycling and actinide production must be modeled.~\cite{cite_07d6df5d6978a1fc,cite_30c88fe46f23448d,cite_612290131318c0d4}

\subsection{Fission Recycling and Abundance Structure}
Fission recycling is the fundamental mechanism governing the nucleosynthetic flow toward the heaviest nuclei. As the r-process path advances along the neutron-drip line, superheavy nuclei undergo spontaneous, beta-delayed, or neutron-induced fission, fragmenting into lighter species that re-enter the process. This recycling stabilizes the abundance distribution, rendering the final pattern largely insensitive to initial seed conditions. However, the precise location of the third r-process peak is contingent upon the timing of neutron release relative to the freeze-out of neutron-capture equilibrium. Late neutron release shifts the peak toward higher mass numbers, while accelerated beta-decay rates can mitigate this shift by promoting earlier fission events. This sensitivity establishes the kinetic delay of fission as a critical, yet under-constrained, variable in abundance predictions.~\cite{cite_b668987a568d271f,cite_fc6f024383b8a4d2}

\subsection{Spectral Diagnostics and Geometric Confounders}
Kilonova spectra provide the primary diagnostic for r-process enrichment, yet they suffer from severe degeneracies between composition and geometry. Lanthanide opacities, driven by dense line forests of f-shell transitions, redshift the emission and prolong cooling times. Viewing-angle-dependent lanthanide curtaining has been observed, but its origin is contested. Geometric asymmetries in the ejecta, such as shock-heated polar winds or tidal tails, and magnetic field configurations can produce anisotropic opacity distributions without invoking time-dependent nuclear kinetics. Consequently, current observational frameworks cannot uniquely attribute spectral evolution to delayed fission recycling; they must first rule out static geometric and magnetic confounders that mimic similar abundance variations.~\cite{cite_43194f3b6113c2ab}

\subsection{The Kinematic Decoupling Gap}
The proposed research addresses a specific unresolved bridge: the lack of a formal criterion distinguishing kinematic fission delays from geometric asymmetries. Existing surveys establish that fission timescales compete with expansion timescales, but they do not provide a unified relation that links the effective fission timescale to both the third r-process peak shift and the temporal evolution of actinide spectral lines. The gap is not merely a deficit in observational data, but a structural absence of a theoretical control panel that can isolate the delayed neutron release mechanism from purely geometric confounders. Without this formal separation, the interpretation of viewing-angle-dependent lanthanide curtaining remains ambiguous.

\subsection{Missing Formal Premises}
The argument for kinematic decoupling requires two specific formal premises that are currently missing from the established literature. First, a precise mathematical definition of the effective fission timescale, incorporating both beta-decay half-lives and isomeric state contributions, is required to compete against the ejecta expansion timescale. Second, a discrimination proof is needed to demonstrate that the temporal evolution of actinide spectral lines under delayed fission is non-isomorphic to patterns generated by static geometric or magnetic asymmetries. These are not empirical observations but necessary theoretical constructs; their absence prevents the formulation of a falsifiable mechanism that can be tested against multi-messenger data.

\subsection{Consequence for the Research Plan}
This gap motivates the shift from empirical surveying to mathematical theory construction. The research plan must therefore establish a static axiomatic framework that defines the kinematic condition where the effective fission timescale exceeds the expansion timescale. This framework will serve as the foundation for deriving the abundance shift and spectral evolution signatures. The transition to the next section will detail the candidate theorem entry, specifying the domain of admissible premises and the scoped entry lemma required to formalize the delayed neutron release mechanism.~\cite{cite_0247e311d8727a8e}

\section{Research Questions and Planned Contributions}
\subsection{Background and Research Gap}
Binary neutron star mergers synthesize heavy elements through the rapid neutron-capture process, where extreme neutron densities drive flow toward the neutron drip line. Fission recycling stabilizes the third r-process peak, but the precise timing of delayed neutron release relative to ejecta expansion remains a critical unresolved variable. The kinematic delay in beta-delayed and isomer-delayed fission neutron release modifies the third r-process peak freeze-out abundance and induces specific temporal evolution of actinide spectral lines. Current models struggle to distinguish whether viewing-angle-dependent lanthanide curtaining arises from kinematic fission delays or purely geometric and magnetic asymmetries in the dynamical ejecta. This ambiguity requires a formal theoretical framework to isolate the temporal signature of delayed fission.~\cite{cite_07d6df5d6978a1fc,cite_30c88fe46f23448d,cite_612290131318c0d4}

\subsection{Operational Research Questions}
\begin{itemize}
\item What is the precise mathematical definition for the 'effective fission timescale' including beta-decay and isomeric state contributions?
\item What is the mathematical formulation of the 'ejecta expansion timescale' in the context of relativistic outflows?
\item What are the operational definitions for the metrics 'third r-process peak abundance shift' and 'spectral evolution'?
\item Can a proof obligation be formulated to demonstrate the non-isomorphism between delayed fission spectral evolution and purely geometric/magnetic asymmetries?
\item What boundary analysis is required to establish the limits of validity for the t\_fission >= t\_expansion assumption?
\item Replace the discarded counterexample\_analyzer batch with qualified human review before treating this design as complete.
\end{itemize}

\subsection{Planned Contributions}
This research plan proposes a unified fission recycling relation that treats the fission timescale as a dynamic variable competing with relativistic expansion. The primary contribution is a formal mechanism linking third r-process peak shifts to unique actinide spectral evolutions and viewing-angle dependencies. By relaxing the instantaneous fission equilibrium assumption, the framework derives an analytical relationship between the effective fission timescale, expansion timescale, and abundance yields. A secondary contribution establishes a discrimination proof obligation to demonstrate the non-isomorphism between spectral evolution under delayed fission versus static geometric or magnetic asymmetry. These formal derivations provide the theoretical basis for interpreting multi-messenger kilonova observations.

\subsection{Design Assumptions and Scope}
The proposed mechanism operates under the design assumption that the effective fission timescale, governed by beta-decay half-lives and nuclear isomeric states, is comparable to or longer than the ejecta expansion timescale in high-entropy, low-density trajectories. Actinide spectral line evolution is assumed to be uniquely sensitive to the fission timescale and not mimicked by variations in ejecta geometry or magnetic field topology alone. The validity domain is bounded by the condition t\_fission >= t\_expansion; the mechanism fails if the effective fission timescale is strictly shorter than the expansion timescale, or if geometric and magnetic asymmetries alone can perfectly replicate the predicted actinide spectral shifts. These assumptions delimit the scope of the formal derivations and the subsequent proof obligations.

\section{Idea Source Checkpoints and Direction Selection Audit}
The selected direction, premise\_inversion, posits that kinematic delay in beta-delayed and isomer-delayed fission neutron release relative to ejecta expansion unifies third r-process peak shifts with viewing-angle-dependent lanthanide curtaining. This mechanism produces a distinct temporal evolution of actinide spectral lines that cannot be replicated by purely geometric or magnetic ejecta asymmetries. The proposal is anchored in the observation that fission recycling dominates the third r-process peak, with fission fragments recycling material into the lanthanide region and influencing the production of lighter heavy elements. The scientific attraction lies in replacing the instantaneous fission equilibrium assumption with a dynamic variable competing with relativistic expansion, thereby linking microphysical nuclear timescales to macroscopic electromagnetic signatures.~\cite{cite_30c88fe46f23448d}

\begin{table*}[htbp]
\centering
\footnotesize
\setlength{\tabcolsep}{3pt}
\renewcommand{\arraystretch}{1.12}
\caption{Decision matrix}
\label{tab:idea-origin-and-selection-checkpoint-table}
\begin{tabular}{|p{0.23\textwidth}|p{0.23\textwidth}|p{0.23\textwidth}|p{0.23\textwidth}|}
\hline
Audit Snapshot & Source File & Direction Mode & Central Hypothesis Focus \\\hline
Checkpoint 1 & idea\_candidate.json & premise\_inversion & Kinematic delay unifying third peak shifts and anisotropic curtaining \\\hline
Checkpoint 2 & idea\_portfolio.json & premise\_inversion & Kinematic delay unifying third peak shifts and spectral signatures \\\hline
Checkpoint 3 & idea\_result.json & premise\_inversion & Kinematic delay unifying third peak shifts and anisotropic curtaining \\\hline
\end{tabular}
\end{table*}

The supplied checkpoints expose critical formal defects that constrain the current theoretical status. The mathematical definitions for the effective fission timescale and the ejecta expansion timescale are absent, preventing the formulation of a rigorous inequality between the two. Furthermore, the specific beta-decay half-life formulas and nuclear isomeric state lifetime models required to compute the fission timescale are not provided. The counterexample analyzer batch was discarded due to an invalid response, leaving the boundary analysis for the t\_fission >= t\_expansion assumption incomplete. Consequently, the proof obligations to demonstrate the analytical relationship between timescales and abundance yields, and the non-isomorphism between delayed fission spectral evolution and geometric asymmetries, remain unverified and require qualified human review.

The retained direction is qualified as a candidate formalization pending the resolution of upstream definitional gaps. The scope is explicitly bounded to high-entropy, low-density trajectories where the effective fission timescale is assumed to be comparable to or longer than the ejecta expansion timescale. The proposal excludes any claim that geometric or magnetic asymmetries cannot produce viewing-angle-dependent lanthanide curtaining; rather, it asserts that they cannot replicate the specific temporal evolution of actinide spectral lines induced by delayed fission. This distinction is central to the premise inversion, which treats the fission timescale as a dynamic variable competing with expansion, while bounding the validity domain against geometric confounders.

\section{Problem Definition, Assumptions, and Hypotheses}
\subsection{Formal Domain and Kinematic Premise}
The formal problem is posed for high-entropy, low-density ejecta trajectories in binary neutron-star mergers, where r-process flow reaches the superheavy region and fission recycling dominates the third abundance peak. Multi-messenger observations establish that kilonova emission is powered by radioactive decay of neutron-rich r-process ejecta, and lanthanide-rich opacity controls the observed redward spectral evolution. The proposed contribution inverts the conventional instantaneous-fission-equilibrium premise: it treats the effective fission timescale as a dynamic variable competing with the ejecta expansion timescale, and asks whether that kinematic delay leaves a unique actinide spectral signature that geometric and magnetic asymmetries cannot reproduce. This section fixes the admissible domain, the premise set, and the failure boundary that the downstream derivation consumes; it does not claim any new observation.~\cite{cite_07d6df5d6978a1fc,cite_612290131318c0d4}

\subsection{Candidate Definition of the Kinematic Ratio}
\paragraph{Definition.} Candidate, unverified definition. Let the kinematic decoupling ratio be the dimensionless quantity R = V1/V2, where V1 denotes the effective fission timescale governed by beta-decay half-lives and nuclear isomeric state lifetimes, and V2 denotes the ejecta expansion timescale of the outflow. The domain of admissible trajectories is the set of high-entropy, low-density ejecta for which fission recycling is active. The instantaneous-fission-equilibrium predicate V10 holds exactly when R is strictly less than one; the delayed-fission regime is entered precisely when R reaches or exceeds one. The formal mathematical definitions of V1 and V2 remain unresolved upstream and are carried as human-input dependencies rather than assumed closed forms.

\subsection{Governing Relation}
R = V1/V2, V10 = (R < 1), delayed-fission regime = (R >= 1)

\subsection{Reading the Governing Relation}
The relation above is the single object the rest of the theory route consumes: the delayed-fission hypothesis is the claim that R >= 1 is physically realized in some admissible trajectories, and that this violation of V10 propagates into a third-peak abundance shift and a distinctive actinide spectral evolution. The equation is a proposed formalization, not a derived identity; its content is fixed only once V1 and V2 receive closed-form definitions. The renderer supplies numbering and labels.

\subsection{Assumption Ledger (Consumer Subset)}
\begin{itemize}
\item This section consumes two premises from the definition ledger and states their scope; the item-by-item ledger is owned elsewhere.
\item A1 (candidate formalization): the effective fission timescale V1, governed by beta-decay half-lives and isomeric state lifetimes, is comparable to or longer than the ejecta expansion timescale V2 in high-entropy, low-density trajectories. This is the delayed-fission premise; it is asserted, not proved, and its validity domain is bounded by the failure condition below.
\item A2 (user-declared): actinide spectral-line evolution V7 is uniquely sensitive to the fission timescale and cannot be mimicked by geometric asymmetries V5 or magnetic-field topology V6 alone. This is a discriminating premise, not an established uniqueness result, and it is exactly what the downstream discrimination obligation must defend.
\end{itemize}

\subsection{Candidate Theorem Entry and Scoped Entry Lemma}
\paragraph{Proof Obligation Registry L1, L2, PO1, PO2 (Candidate).} Candidate theorem entry (proposed, unverified). Candidate L1: under premises A1 and A2, the violation of the instantaneous-equilibrium predicate (R >= 1) yields a third r-process peak abundance shift that drives a temporal evolution of actinide spectral lines statistically distinct from patterns generated by V5 or V6 alone. Candidate L2: under A2, no mapping from geometric or magnetic asymmetry configurations reproduces the delayed-fission spectral evolution. These are scoped entry lemmas, not proved results; they package the two hypotheses P1 and P2 as the domain-and-boundary statement the derivation stage inherits. The associated proof obligations PO1 (an analytical relation among R, V1, V2, and the abundance shift) and PO2 (a non-isomorphism proof for the spectral signature) are open and require human confirmation of the unresolved definitions before either lemma can be discharged.

\subsection{Failure Boundary and Transition}
The mechanism fails on either of two declared boundaries: if the effective fission timescale for superheavy nuclei is strictly shorter than the expansion timescale (R < 1 throughout the admissible domain), or if geometric and magnetic asymmetries alone can perfectly replicate the predicted actinide spectral shifts without invoking a delayed fission timescale. These boundaries are scope delimiters, not evidence of failure; no observation is reported, and the frozen counterexample analysis was discarded pending human review. The next stage, the definition and proposition ledger, supplies the closed-form definitions of V1, V2, V3, V7 and the derivation lemmas that this entry consumes.

\section{Study Design and Methods}
\subsection{Static Axiomatic Framework and Kinematic Regimes}
This study adopts a static axiomatic framework to formalize the kinematic decoupling of fission recycling in kilonovae. The design centers on the competition between the effective fission timescale and the ejecta expansion timescale. Within this theoretical structure, the mechanism operates under a specific regime where the effective fission timescale is comparable to or longer than the ejecta expansion timescale. This condition violates the instantaneous equilibrium predicate, which assumes that the fission timescale is strictly shorter than the expansion timescale. By relaxing the instantaneous fission equilibrium assumption and substituting it with a delayed beta and isomer fission release axiom, the framework treats the fission timescale as a dynamic variable competing with relativistic expansion. This approach directly links the modification of the third r-process peak freeze-out abundance to viewing-angle-dependent lanthanide curtaining, providing a formal mechanism that connects microphysical nuclear delays to macroscopic electromagnetic signatures.

\subsection{Comparators and Boundary Conditions}
The proposed theoretical contribution is evaluated against specific comparators and strict boundary conditions. The primary control variable is the instantaneous fission equilibrium limit, which serves as the baseline against which delayed fission effects are measured. The design explicitly accounts for confounders, namely geometric asymmetries and magnetic field configurations in the ejecta. The mechanism is bounded by the assumption that actinide spectral line evolution is uniquely sensitive to the fission timescale. Consequently, the theory fails if geometric or magnetic asymmetries alone can perfectly replicate the predicted actinide spectral shifts without invoking delayed fission timescales. This boundary condition ensures that the proposed relation is not conflated with alternative explanations that attribute viewing-angle-dependent lanthanide curtaining to purely structural or magnetic effects in the dynamical ejecta.

\subsection{Protocol and Dependency Management}
The protocol for this mathematical theory route relies on the derivation of analytical relationships and the demonstration of non-isomorphism between distinct spectral evolution patterns. The derivation requires formal definitions for the effective fission timescale, the ejecta expansion timescale, and the specific metrics used to quantify abundance shifts and spectral evolution. These operational definitions are currently unresolved and depend on upstream human input to finalize the mathematical formulations. Furthermore, the quality control and measurement plans required to translate these theoretical constructs into observational diagnostics are pending human confirmation. The study design integrates these dependencies by specifying that the formal reasoning plan and counterexample analysis require qualified human review before the design can be considered complete. This ensures that the theoretical framework remains rigorous and that all mathematical premises are explicitly defined before further derivation steps are executed.

\subsection{Expected Outcomes and Decision Rules}
The expected outcomes of this theoretical design are mapped to specific decision rules that govern the status of the proposed mechanisms. If the prespecified analysis is consistent with the declared relation and the planned controls do not favor the alternative explanation of pure geometric or magnetic asymmetry, the result supports the mechanism within the defined boundary. Conversely, if the comparison does not support the declared relation or favors the alternative explanation, the result is null or contradictory, indicating that the proposed relation is not supported in this design boundary. A partial or heterogeneous outcome occurs if the analysis indicates variation across declared conditions, units, or measurement contexts, suggesting that the relation may be conditional. Finally, an uninformative or invalid outcome arises if prespecified quality-control, missingness, or validity criteria prevent interpretation, necessitating the resolution of identified validity or data-quality failures before repeating the design.

\subsection{Data Governance and Reproducibility}
The data governance and reproducibility protocols for this theoretical framework require human input to finalize the management and reproducibility standards. Because the study design is a mathematical theory route, the primary artifacts are formal proofs, derivations, and analytical relationships rather than empirical datasets. However, the reproducibility of these formal proofs depends on the precise operational definitions of the variables involved, which are currently pending resolution. The data management plan must ensure that all upstream formal records, outcome branches, and idea checkpoints are globally available and consistently referenced. This governance structure supports the auditability of the theoretical spine, ensuring that every lemma, proof obligation, and decision branch can be traced back to its foundational premises and source materials.

\section{Expected Outcome Branches and Conditional Conclusions}
\subsection{Decision Frame for the Kinematic Decoupling Proposal}
This section converts the proposed delayed-fission mechanism into a preregistered decision protocol. The upstream derivation supplies the controlling premises: the effective fission timescale is assumed to compete with the ejecta expansion timescale, and the actinide spectral signature is declared non-isomorphic to purely geometric or magnetic asymmetries. The protocol maps each prespecified outcome branch to the relevant lemma and proof obligation, the conclusion permitted within the declared boundary, and the next action. No branch is treated as observed; all consequences are conditional on the formal comparison specified by the design.

\begin{table*}[htbp]
\centering
\footnotesize
\setlength{\tabcolsep}{3pt}
\renewcommand{\arraystretch}{1.12}
\caption{Outcome-to-Lemma-to-Action Matrix}
\label{tab:expected-outcomes-decision-matrix}
\begin{tabular}{|p{0.184\textwidth}|p{0.184\textwidth}|p{0.184\textwidth}|p{0.184\textwidth}|p{0.184\textwidth}|}
\hline
Outcome branch & Trigger & Relevant Lemma / PO & Allowed conclusion & Next action \\\hline
supports\_mechanism & Analysis consistent with declared relation; controls do not favor geometric or magnetic alternatives & L1, PO1, PO2 & Relation supported within design boundary; no proof of the theorem & Replicate under independently confirmed conditions; test the most consequential declared boundary condition \\\hline
partial\_or\_heterogeneous & Variation across declared conditions, units, or measurement contexts & L2, PO2 & Conditional or heterogeneous relation; no universal conclusion & Predefine and check plausible moderators; improve coverage and measurement comparability \\\hline
null\_or\_contradictory & Comparison does not support the relation or favors a declared alternative & L3, PO1 & Relation not supported in this boundary; absence of support is not general proof of absence & Audit construct validity and comparison adequacy; revise mechanism or boundary claim \\\hline
uninformative\_or\_invalid & Quality-control, missingness, protocol-deviation, or validity criteria prevent interpretation & L4, L5, L6, PO1, PO2 & No scientific conclusion warranted; theorem status does not update & Resolve validity or data-quality failure; obtain human confirmation of prerequisites \\\hline
\end{tabular}
\end{table*}

\subsection{Supportive Branch and Its Scope}
\paragraph{Pre-registered Branch (Expected---Not Observed).} The supportive branch is reached only when the prespecified analysis is consistent with the declared relation while the planned controls fail to favor a geometric or magnetic alternative. Under that trigger, the permitted conclusion is support, not proof: the delayed-fission mechanism remains a candidate relation confined to the design boundary. The next actions are replication under independently confirmed conditions and a direct test of the most consequential declared boundary condition, namely the comparison of the effective fission timescale with the expansion timescale.

\subsection{Heterogeneous and Competing-Explanation Branches}
\paragraph{Pre-registered Branch (Expected---Not Observed).} The heterogeneous branch records variation across declared conditions, units, or measurement contexts; it licenses only a conditional reading and directs the work toward predefined moderators and improved comparability. The null or contradictory branch records a comparison that fails to support the relation or instead favors the declared alternative explanations, in which lanthanide curtaining and abundance variation are attributed to geometric asymmetries or magnetic field configurations rather than kinematic fission delays. That outcome does not refute the mechanism generally; it confines the proposal to the tested boundary and triggers an audit of construct validity and comparison adequacy before any revision.

\subsection{Uninformable Branch and Dependency Consequence}
\paragraph{Pre-registered Branch (Expected---Not Observed).} The uninformative or invalid branch is entered when prespecified quality-control, missingness, protocol-deviation, or validity criteria prevent interpretation. This branch is a no-information condition for the dependency: it neither confirms nor denies the proposed relation, and it does not update theorem status. The required response is to resolve the identified validity or data-quality failure and to obtain human confirmation of the measurement, sampling, and analysis prerequisites before repeating the design. Treating this branch as a proof of failure would misread the protocol; it is a procedural dependency, not a scientific result.

\subsection{Transition to Risk and Review}
The matrix above fixes the conditional conclusions and next actions for each branch. The detailed item-by-item human-confirmation and release criteria are owned by the review ledger and are not restated here. The plan now carries the scoped dependency on that ledger forward: the decision protocol is executable only once the review gate confirms the prerequisites that the uninformative branch protects against.

\section{Risks, Limitations, and Human Review Requirements}
\begin{table*}[htbp]
\centering
\footnotesize
\setlength{\tabcolsep}{3pt}
\renewcommand{\arraystretch}{1.12}
\caption{Release and Review Decision Matrix}
\label{tab:risk-limitations-and-review-risk-limitations-and-review-1}
\begin{tabular}{|p{0.184\textwidth}|p{0.184\textwidth}|p{0.184\textwidth}|p{0.184\textwidth}|p{0.184\textwidth}|}
\hline
Review Item & Scope & Required Human Decision & Release Condition & Consequence for Theory Status \\\hline
Counterexample analysis & Formal-theory boundary testing & Qualified expert review of the discarded counterexample batch & Re-run or human adjudication of candidate counterexamples against all declared assumptions & No conclusion about counterexamples; the mechanism remains unverified \\\hline
Formal reasoning plan & Lemma, definition, and proof-obligation ledger & Review of the unverified formal derivation chain & Approval of definitions and proof obligations before any theorem-status update & The candidate theorem remains a proposal, not a verified result \\\hline
Risk gate & Overall release authorization & Confirmation that no empirical claim is embedded in the formal plan & Sign-off that the proposal is scoped to mathematical theory only & The plan may proceed as a design artifact, not as an observed finding \\\hline
Workflow degradation & LLM or workflow integrity & Inspection of the discarded analyzer batch and its downstream effects & Documentation that the discarded batch introduced no retained claim or protocol detail & The degraded stage is treated as a procedural dependency, not a scientific result \\\hline
\end{tabular}
\end{table*}

\subsection{Scope Limits and Alternative Mechanisms}
The proposal's central limitation is the absence of a completed formal derivation and the corresponding lack of observed empirical validation. The frozen evidence establishes that the kinematic decoupling of fission recycling is a candidate mechanism, not a demonstrated one. Two competing explanations remain live: geometric asymmetries in dynamical ejecta, such as shock-heated polar winds and tidal tails, and magnetic field configurations. Both can produce viewing-angle-dependent lanthanide curtaining without invoking delayed fission timescales. The design distinguishes the proposed mechanism from these alternatives only through the temporal evolution of actinide spectral lines, a signature that the current plan has not yet formalized or observed. The control condition, instantaneous fission equilibrium, is declared but not analytically bounded. Consequently, the comparison between delayed fission and static geometric or magnetic asymmetries is a design requirement, not an established result. The proposal's scope is confined to the high-entropy, low-density trajectories where the effective fission timescale is assumed to be comparable to or longer than the ejecta expansion timescale; outside that domain, the mechanism is not addressed.

\subsection{Procedural Dependencies and No-Information Branches}
\paragraph{Decision Status: No-information.} The theory spine records no-information branches for the unverified proof obligations and for the undefined operational metrics. These branches do not indicate a failed derivation; they indicate that the dependency does not update theorem status. The effective fission timescale, the ejecta expansion timescale, the abundance-shift metric, and the spectral-evolution metric each require a formal definition before the candidate lemmas can be evaluated. Until those definitions are supplied, the proof obligations remain open and the lemmas remain candidate. The discarded counterexample batch is a procedural gap, not a scientific finding: no counterexample is retained, and no conclusion about the mechanism's validity is drawn from its absence. The review decision matrix above maps each of these gaps to a concrete human action. The transition to the next stage is therefore conditional: the mathematical-theory route proceeds only after the definition ledger and the proof-obligation registry are closed by qualified review, at which point the derivation and falsification sections can consume a settled premise set rather than a placeholder set.

\subsection{Concrete Release Conditions}
\paragraph{Human-review checklist.} 1. Confirm that the counterexample analysis is either re-executed or formally waived by a qualified reviewer, and that no conclusion about the mechanism's validity is drawn from the discarded batch. 2. Approve the formal reasoning plan's definitions for the effective fission timescale, the ejecta expansion timescale, the abundance-shift metric, and the spectral-evolution metric; without these, the candidate lemmas cannot be evaluated. 3. Verify that the risk gate is scoped to mathematical theory only, with no empirical claim embedded in the formal derivation chain. 4. Document that the workflow degradation introduced no retained claim, formalization, counterexample, or protocol detail, and that the degraded stage is treated as a procedural dependency. 5. Authorize the next stage only after the definition ledger and proof-obligation registry are closed, so that the derivation and falsification sections consume settled premises rather than placeholders.

\section{Definitions, Propositions, and Proof Obligations}
\begin{table}[htbp]
\centering
\small


\caption{Definition Ledger and Symbol Domain}
\label{tab:definitions-and-propositions-def-ledger}
\begin{tabular}{|p{0.29\linewidth}|p{0.29\linewidth}|p{0.29\linewidth}|}
\hline
Symbol & Role & Status \\\hline
V1 & Effective fission timescale & needs\_human\_input \\\hline
V2 & Ejecta expansion timescale & needs\_human\_input \\\hline
V3 & Third r-process peak abundance & needs\_human\_input \\\hline
V4 & Lanthanide curtaining metric & needs\_human\_input \\\hline
V5 & Geometric asymmetries & needs\_human\_input \\\hline
V6 & Magnetic field configurations & needs\_human\_input \\\hline
V7 & Actinide spectral evolution & needs\_human\_input \\\hline
V8 & Beta-decay half-lives & needs\_human\_input \\\hline
V9 & Isomeric state lifetimes & needs\_human\_input \\\hline
V10 & Instantaneous equilibrium & candidate\_formalization \\\hline
\end{tabular}
\end{table}

\subsection{Assumption Ledger}
\begin{itemize}
\item A1: The effective fission timescale is comparable to or longer than the ejecta expansion timescale in high-entropy, low-density trajectories (candidate\_formalization).
\item A2: Actinide spectral line evolution is uniquely sensitive to the fission timescale and cannot be mimicked by variations in ejecta geometry or magnetic field topology alone (user\_declared).
\end{itemize}

\subsection{Kinematic Threshold Relation}
\begin{equation}
V1 \geq V2
\label{eq:definitions-and-propositions-eq-threshold}
\end{equation}

\subsection{Proposition P1: Delayed Neutron Release Mechanism}
\paragraph{Proposition (Candidate).} Candidate Proposition: Delayed neutron release due to finite fission timescale (V1 >= V2) modifies the third r-process peak abundance (V3) and induces specific temporal evolution of actinide spectral lines (V7). See Eq.~\eqref{eq:definitions-and-propositions-eq-threshold}.

\subsection{Proposition P2: Spectral Discrimination}
\paragraph{Proposition (Candidate).} Candidate Proposition: The specific temporal evolution of actinide spectral lines (V7) caused by delayed fission cannot be replicated by geometric (V5) or magnetic (V6) asymmetries alone. See Eq.~\eqref{eq:definitions-and-propositions-eq-threshold}.

\begin{table*}[htbp]
\centering
\footnotesize
\setlength{\tabcolsep}{3pt}
\renewcommand{\arraystretch}{1.12}
\caption{Proof Obligation Registry and Dependency Matrix}
\label{tab:definitions-and-propositions-po-registry}
\begin{tabular}{|p{0.23\textwidth}|p{0.23\textwidth}|p{0.23\textwidth}|p{0.23\textwidth}|}
\hline
Obligation & Scope & Required Inputs & Status \\\hline
PO1 & Kinematic decoupling to abundance shift & A1, D1, D2, D3, P1, S1, S2, S3 & unverified \\\hline
PO2 & Spectral non-isomorphism to asymmetries & A2, D5, D8, D9, P1, P2, S3, S4 & unverified \\\hline
\end{tabular}
\end{table*}

\subsection{Failure Conditions and Unresolved Inputs}
\paragraph{Decision Status: No-information.} The mechanism fails if the effective fission timescale for superheavy nuclei is strictly shorter than the expansion timescale, or if geometric/magnetic asymmetries alone can perfectly replicate the predicted actinide spectral shifts without invoking delayed fission timescales. The formalization currently requires human input to define the precise mathematical formulations for the effective fission timescale (V1), ejecta expansion timescale (V2), beta-decay half-lives (V8), and isomeric state lifetimes (V9). These missing definitions constitute no-information branches that withhold theorem status updates.

\section{Forward Derivation and Counterexample Search Plan}
\subsection{Derivation Domain and Admissible Premises}
\paragraph{Proof Obligation Registry PO1, PO2 (Unverified).} This section advances the argument by specifying the derivation obligations and counterexample search plan for the kinematic decoupling mechanism. The formal domain is restricted to high-entropy, low-density ejecta trajectories where the effective fission timescale V1 and the ejecta expansion timescale V2 are the primary independent variables. The admissible premises are the candidate formalization A1 and the user-declared A2. The derivation does not assume the truth of the definitions for V1, V2, V3, V7, V8, or V9; rather, it treats them as placeholders that must be resolved by the definitions ledger to close the proof obligations. The argument proceeds by establishing the logical consequences of the timescale inequality V1 >= V2 on the third r-process peak abundance V3 and the actinide spectral evolution V7, subject to the constraint that geometric V5 and magnetic V6 asymmetries are controlled confounders.

\subsection{Kinematic Decoupling Condition}
\begin{equation}
V1 >= V2
\label{eq:forward-derivation-and-counterexamples-b2}
\end{equation}

\subsection{Derivation Chain and Lemma Registry}
The derivation chain proceeds through four candidate lemmas. Lemma L3 applies the kinematic condition to violate the instantaneous equilibrium predicate, establishing the regime where delayed neutron release is dynamically active. Lemma L4 propagates this violation to the nucleosynthesis yield, asserting that the shift in the third r-process peak abundance V3 is a direct consequence of the finite beta-decay half-lives V8 and isomeric state lifetimes V9. Lemma L5 connects the abundance shift to the spectral domain, proposing that the temporal evolution of actinide lines V7 is statistically distinct from patterns generated by geometric V5 or magnetic V6 asymmetries. Finally, Lemma L6 unifies the viewing-angle-dependent lanthanide curtaining V4 with the spectral evolution V7, closing the argument that delayed fission recycling provides a coherent mechanism for the observed phenomena. Each lemma is unverified and depends on the resolution of the upstream definitions.

\subsection{Candidate Derivation Obligations}
\paragraph{Proof Obligation Registry PO1, PO2 (Unverified).} Two proof obligations must be discharged to validate the derivation. Proof Obligation PO1 requires the derivation of an analytical relationship between the effective fission timescale V1, the expansion timescale V2, and the abundance yields V3, conditional on the definitions of beta-decay half-lives V8 and isomeric lifetimes V9. Proof Obligation PO2 requires a discrimination proof demonstrating the non-isomorphism between the spectral evolution V7 under delayed fission and the static patterns produced by geometric V5 or magnetic V6 asymmetries alone. These obligations are currently unverified and require human input to formalize the mathematical definitions of the variables.

\begin{table*}[htbp]
\centering
\footnotesize
\setlength{\tabcolsep}{3pt}
\renewcommand{\arraystretch}{1.12}
\caption{Counterexample Decision Matrix}
\label{tab:forward-derivation-and-counterexamples-b5}
\begin{tabular}{|p{0.23\textwidth}|p{0.23\textwidth}|p{0.23\textwidth}|p{0.23\textwidth}|}
\hline
Decision & Condition & Interpretation & Action \\\hline
Valid Counterexample & V1 < V2 in all high-entropy trajectories & Mechanism fails at the kinematic level & Revise assumption A1 \\\hline
Valid Counterexample & V7 is isomorphic to V5 or V6 patterns & Mechanism fails at the spectral discrimination level & Reject proposition P2 \\\hline
Boundary Case & V1 approx V2 with high variance & Mechanism is conditional on specific trajectory details & Define variance bounds for V1 \\\hline
Out-of-Domain & Low-entropy or high-density ejecta & Premises A1 and A2 do not apply & Exclude from derivation scope \\\hline
No Information & Definitions for V1, V2, or V7 are missing & Logical status of lemmas cannot be evaluated & Resolve definition ledger entries \\\hline
\end{tabular}
\end{table*}

\subsection{Failure Conditions and Limitations}
\begin{itemize}
\item The derivation is subject to strict failure conditions. First, if the effective fission timescale V1 is strictly shorter than the ejecta expansion timescale V2, the mechanism fails to produce the predicted abundance shifts. Second, if geometric asymmetries V5 or magnetic field configurations V6 alone can perfectly replicate the predicted actinide spectral shifts V7 without invoking delayed fission timescales, the uniqueness claim of the mechanism fails. These conditions define the boundary of the proposal's validity. The counterexample analysis batch was discarded, and no specific counterexamples have been generated; therefore, the search plan remains a procedural dependency requiring qualified human review to complete the verification cycle.
\end{itemize}

\appendices
\section{Energy-Condition Taxonomy, Symbols, and Boundary Defense}
\subsection{Context and Scope of the Boundary Defense}
This appendix establishes the logical boundaries for the proposed kinematic decoupling mechanism in neutron star merger nucleosynthesis. The central hypothesis posits that the effective fission timescale (\$V\_1\$) is comparable to or longer than the ejecta expansion timescale (\$V\_2\$), violating instantaneous equilibrium and driving third r-process peak shifts and unique actinide spectral evolutions. This defense clarifies that the formal variables \$V\_1\$ through \$V\_\{10\}\$ are operational constructs for nucleosynthetic kinetics and radiative transfer, distinct from the geometric energy conditions (NEC, SEC, AANEC) governing spacetime curvature. Consequently, no implication from general relativistic focusing theorems is assumed or required for the validity of the proposed mechanism.

\begin{table*}[htbp]
\centering
\footnotesize
\setlength{\tabcolsep}{3pt}
\renewcommand{\arraystretch}{1.12}
\caption{Taxonomy of Energy Conditions and Kinematic Boundaries}
\label{tab:appendix-variables-and-definitions-B2}
\begin{tabular}{|p{0.3067\textwidth}|p{0.3067\textwidth}|p{0.3067\textwidth}|}
\hline
Condition & Domain & Logical Relation to Proposed Mechanism \\\hline
**NEC** & Spacetime Geometry & Independent. Governs null energy flux; does not constrain nuclear decay rates or ejecta opacity. \\\hline
**AANEC** & Spacetime Geometry & Independent. Averaged null energy condition; establishes focusing in GR but does not imply kinematic decoupling of fission. \\\hline
**SEC** & Spacetime Geometry & Independent. Strong energy condition; irrelevant to the timescale competition between \$V\_1\$ and \$V\_2\$. \\\hline
**\$V\_1 \textbackslash{}ge V\_2\$** & Nuclear Kinetics & Proposed Mechanism. Effective fission timescale exceeds expansion timescale, violating instantaneous equilibrium predicate \$V\_\{10\}\$. \\\hline
**\$V\_7\$ Sensitivity** & Spectral Physics & Proposed Discriminator. Actinide line evolution is uniquely sensitive to \$V\_1\$ and cannot be mimicked by geometry (\$V\_5\$) or magnetic fields (\$V\_6\$). \\\hline
\end{tabular}
\end{table*}

\subsection{Defense Against Automatic Interchangeability}
The stress-energy tensor of the ejecta does not automatically satisfy or violate geometric energy conditions in a way that dictates nucleosynthetic outcomes. Specifically, the satisfaction of the Null Energy Condition (NEC) by the matter content does not establish the kinematic delay \$V\_1 \textbackslash{}ge V\_2\$; conversely, the failure of the Strong Energy Condition (SEC) does not preclude fission recycling. The proposed mechanism relies entirely on the competition between nuclear weak-interaction rates and hydrodynamic expansion, as formalized in the candidate lemma \$L\_3\$. AANEC, while relevant for trapped surface formation in extreme gravity, provides no information regarding the temporal evolution of actinide spectral lines (\$V\_7\$). Therefore, the boundary between geometric constraints and kinetic modeling is strictly maintained: no geometric condition substitutes for the explicit assumption of delayed beta-and-isomer fission release.

\subsection{Operational Definitions and Human Review Dependencies}
The formalization of the mechanism is contingent upon the precise mathematical definition of the effective fission timescale (\$V\_1\$) and the ejecta expansion timescale (\$V\_2\$), which currently require human review to specify thresholds for 'comparable' versus 'longer'. Additionally, the quantitative metrics for third r-process peak freeze-out abundance (\$V\_3\$) and viewing-angle-dependent lanthanide curtaining (\$V\_4\$) are undefined in the source material. The spectral evolution metric (\$V\_7\$) lacks specification of the exact actinide lines (e.g., Uranium or Thorium) and their corresponding wavelengths. These unresolved operational definitions constitute the primary procedural dependencies for advancing the proof obligations \$PO\_1\$ and \$PO\_2\$.

\section{Idea Source Checkpoints and Direction Selection Audit}
\subsection{Premise Inversion and Mechanistic Unification}
The selected direction, premise inversion, replaces the instantaneous fission equilibrium assumption with a delayed beta and isomer fission release axiom. This shift treats the effective fission timescale as a dynamic variable competing with the ejecta expansion timescale. The mechanism unifies third r-process peak shifts with viewing-angle-dependent lanthanide curtaining, producing a distinct temporal evolution of actinide spectral lines. This theoretical pivot is grounded in established nuclear fission dynamics and r-process network calculations, which demonstrate that fission recycling governs the nucleosynthetic flow toward the heaviest nuclei. The kinetic delay in neutron release relative to expansion modifies the freeze-out abundance, a relation supported by the fundamental competition between neutron capture and beta decay in extreme astrophysical environments.~\cite{cite_30c88fe46f23448d,cite_b668987a568d271f,cite_fc6f024383b8a4d2}

\subsection{Spectral Discrimination Against Geometric Confounders}
The central hypothesis asserts that the temporal evolution of actinide spectral lines is uniquely sensitive to the fission timescale and cannot be mimicked by variations in ejecta geometry or magnetic field topology alone. This discriminative claim is critical for isolating the kinematic delay mechanism from alternative explanations. Multi-messenger observations of neutron star mergers reveal distinct high-opacity lanthanide-rich ejecta and lower-opacity components, where viewing-angle dependencies and asymmetries in ejecta morphology significantly alter the observed spectral energy distribution. By bounding the validity domain against these geometric confounders, the proposed framework requires a spectroscopic identification of actinide line evolution that is strictly inconsistent with both instantaneous fission equilibrium and pure geometric asymmetry models.~\cite{cite_07d6df5d6978a1fc,cite_43194f3b6113c2ab,cite_612290131318c0d4}

\begin{table*}[htbp]
\centering
\footnotesize
\setlength{\tabcolsep}{3pt}
\renewcommand{\arraystretch}{1.12}
\caption{Audit Snapshots and Retained Decisions}
\label{tab:appendix-idea-evolution-b3}
\begin{tabular}{|p{0.3067\textwidth}|p{0.3067\textwidth}|p{0.3067\textwidth}|}
\hline
Checkpoint & Source File & Retained Decision \\\hline
Audit Snapshot 1 & idea\_candidate.json & Establishes the core premise inversion and unified fission recycling relation. \\\hline
Audit Snapshot 2 & idea\_portfolio.json & Refines the title to emphasize spectral signatures alongside third r-process peak shifts. \\\hline
Audit Snapshot 3 & idea\_result.json & Details the boundary conditions, falsifiers, and the requirement to isolate the actinide temporal signature from geometric confounders. \\\hline
\end{tabular}
\end{table*}

\section{Evidence Coverage, Unknown Items, and Review Checklist}
\subsection{Evidence Coverage and Boundary Use}
The proposal's physical premise is anchored in established multi-messenger evidence: binary neutron star mergers synthesize heavy r-process elements whose lanthanide and actinide content governs kilonova opacity and spectral evolution (W2766840380; W27655595587). Dynamical ejecta with low electron fraction drive fission recycling, which shapes the third abundance peak and produces viewing-angle-dependent lanthanide curtaining (W2760756222; W3092825069). These sources establish the astrophysical environment in which delayed fission neutron release operates, but they do not demonstrate the proposed kinematic decoupling mechanism. The evidence supports the phenomenon and the relevant physical scales; it does not validate the candidate theorem linking fission timescale to spectral discrimination. The contribution of this route is therefore a formal design that isolates the delayed-fission signature from geometric and magnetic confounders, not an empirical confirmation of that signature.~\cite{cite_07d6df5d6978a1fc,cite_612290131318c0d4,cite_0247e311d8727a8e}

\subsection{Unknown-Item Consequences for the Theory Spine}
The degraded counterexample-analyzer stage (unknown-8) leaves the proposal without a machine-generated counterexample ledger, so no conclusion about the existence or non-existence of counterexamples may be drawn. This gap is consequential because the central claim depends on a non-isomorphism condition: that delayed-fission spectral evolution cannot be replicated by geometric or magnetic asymmetries alone. Without a completed counterexample analysis, that condition remains an unverified assumption rather than a defended boundary. The theory spine inherits this dependency: the proof obligations for the abundance-shift lemma and the spectral-discrimination lemma cannot be discharged until the counterexample space is bounded by qualified human review. The absence of results is not evidence of failure, but it is a precise procedural dependency that blocks theorem-status updates.

\begin{table*}[htbp]
\centering
\footnotesize
\setlength{\tabcolsep}{3pt}
\renewcommand{\arraystretch}{1.12}
\caption{Release Criteria for Future Claims}
\label{tab:appendix-evidence-and-review-b3}
\begin{tabular}{|p{0.3067\textwidth}|p{0.3067\textwidth}|p{0.3067\textwidth}|}
\hline
Claim Type & Release Criterion & Blocking Dependency \\\hline
Candidate abundance-shift lemma & Formal definitions of effective fission timescale and expansion timescale supplied and internally consistent & Human review of definition ledger entries D1, D2 \\\hline
Candidate spectral-discrimination lemma & Counterexample space bounded; non-isomorphism condition tested against declared confounders & Qualified human review of degraded counterexample-analyzer stage \\\hline
Outcome-branch assignment & Prespecified analysis executed with controls for geometric and magnetic asymmetries & Human confirmation of measurement and analysis prerequisites \\\hline
\end{tabular}
\end{table*}\bibliographystyle{IEEEtran}
\bibliography{references}
\end{document}
